% Reproducible build:
%   tectonic docs/PLUGIN_STANDARD.tex --outdir /tmp/plugin-standard-build
%   install -m 0644 /tmp/plugin-standard-build/PLUGIN_STANDARD.pdf \
%     assets/documents/vllm-hust-plugin-standard-v1.0.pdf
\documentclass[11pt,a4paper,fontset=none]{ctexart}

% Legacy compatibility profile. The organization-wide ecosystem architecture
% now separates system role, integration contract, execution plane, deployment
% topology, and delivery model. This document only covers the existing Python
% entry-point delivery mechanism during migration.

\usepackage[a4paper,margin=22mm,headheight=16pt]{geometry}
\usepackage{array}
\usepackage{amssymb}
\usepackage{booktabs}
\usepackage{enumitem}
\usepackage{fancyhdr}
\usepackage{hyperref}
\usepackage{listings}
\usepackage{longtable}
\usepackage{microtype}
\usepackage[most]{tcolorbox}
\usepackage{xcolor}

\setCJKmainfont{FandolSong-Regular.otf}[BoldFont=FandolSong-Bold.otf]
\setCJKsansfont{FandolSong-Regular.otf}[BoldFont=FandolSong-Bold.otf]
\setCJKmonofont{FandolSong-Regular.otf}

\definecolor{Ink}{HTML}{0C1112}
\definecolor{Teal}{HTML}{176F72}
\definecolor{Cyan}{HTML}{39C6CA}
\definecolor{Lime}{HTML}{A8D951}
\definecolor{Mist}{HTML}{E9EFEC}
\definecolor{Slate}{HTML}{52615F}
\definecolor{Amber}{HTML}{F4E1AE}

\hypersetup{
  colorlinks=true,
  linkcolor=Teal,
  urlcolor=Teal,
  pdftitle={vLLM-HUST Plugin Standard 1.0},
  pdfauthor={vLLM-HUST},
  pdfsubject={Development and operations contract for vLLM runtime plugins},
  pdfkeywords={vLLM, plugin, inference engine, lifecycle, conformance}
}
\Urlmuskip=0mu plus 2mu

\setlength{\parindent}{0pt}
\setlength{\parskip}{0.55em}
\setlist{nosep,leftmargin=2em}
\renewcommand{\arraystretch}{1.25}
\pagestyle{fancy}
\fancyhf{}
\lhead{\textcolor{Teal}{vLLM-HUST Plugin Standard}}
\rhead{\textcolor{Slate}{Version 1.0}}
\cfoot{\textcolor{Slate}{\thepage}}
\renewcommand{\headrulewidth}{0.4pt}

\lstdefinestyle{standardcode}{
  basicstyle=\ttfamily\small,
  backgroundcolor=\color{Mist},
  frame=single,
  rulecolor=\color{Teal!35},
  breaklines=true,
  columns=fullflexible,
  keepspaces=true,
  showstringspaces=false,
  aboveskip=0.8em,
  belowskip=0.8em
}
\lstset{style=standardcode}

\newtcolorbox{normative}[1]{
  colback=Mist,
  colframe=Teal,
  title=\textbf{#1},
  fonttitle=\bfseries,
  boxrule=0.8pt,
  arc=0pt
}
\newcommand{\must}{\textbf{MUST / 必须}}
\newcommand{\should}{\textbf{SHOULD / 应当}}
\newcommand{\may}{\textbf{MAY / 可以}}

\begin{document}

\begin{titlepage}
  \pagecolor{Ink}\color{white}
  \vspace*{20mm}
  {\color{Cyan}\bfseries\sffamily INFERENCE ENGINE EXTENSION CONTRACT\par}
  \vspace{15mm}
  {\fontsize{34}{40}\selectfont\bfseries\sffamily vLLM-HUST\par}
  \vspace{3mm}
  {\fontsize{27}{34}\selectfont\bfseries\sffamily Plugin Standard 1.0\par}
  \vspace{5mm}
  {\Large 插件开发与运行标准\par}
  \vfill
  \begin{tcolorbox}[colback=white!7!Ink,colframe=white!25,arc=0pt,boxrule=0.6pt]
    \color{white}
    本标准规定 vLLM-HUST 运行时插件的最小开发、集成、启动、验证、停止、卸载和回滚合同。\\
    This standard defines the minimum development, integration, startup, verification,
    shutdown, removal, and rollback contract for a vLLM-HUST runtime plugin.
  \end{tcolorbox}
  \vspace{12mm}
  {\color{Lime}\bfseries Version 1.0}\hfill 2026-08-26
  \vspace{8mm}
\end{titlepage}
\nopagecolor\color{Ink}

\section*{Document status / 文档状态}
\addcontentsline{toc}{section}{Document status / 文档状态}

This is a stable public engineering standard. Conformance refers to a specific plugin release and
its declared compatibility range; it does not imply performance superiority or endorsement of a
research claim.

本文件是稳定的公开工程标准。“符合本标准”仅针对某一插件版本及其声明的兼容范围，不代表性能
优越，也不构成对研究结论的背书。

\begin{longtable}{>{\bfseries}p{0.28\textwidth}p{0.64\textwidth}}
Version / 版本 & 1.0 \\
Status / 状态 & Public standard / 公开标准 \\
Applies to / 适用对象 & Independently packaged vLLM runtime plugins / 独立封装的 vLLM 运行时插件 \\
Canonical source / 权威源 & \url{https://github.com/vLLM-HUST/vllm-hust-website/blob/main/docs/PLUGIN_STANDARD.tex} \\
Project site / 项目网站 & \url{https://vllm-hust.sage.org.ai/plugins.html} \\
\end{longtable}

\tableofcontents
\newpage

\section{Normative language / 规范性语言}

The key words \must, \should, and \may indicate requirement, recommendation, and permission.
Where English and Chinese wording differ, the requirement with the stricter operational safety
interpretation applies until the source is corrected.

关键词 \must、\should、\may 分别表示强制要求、建议要求和允许事项。中英文出现差异时，在源文件
修订前采用运行安全约束更严格的解释。

\section{Scope and boundary / 范围与边界}

A conforming plugin is an independently versioned Python distribution that extends vLLM through a
supported Python entry-point group. It installs without modifying the vLLM-HUST source tree and has
a documented enable, verify, disable, remove, and rollback path.

合规插件是独立版本化的 Python 发行包，通过受支持的 Python entry point 扩展 vLLM。安装插件不得
修改 vLLM-HUST 源码树，并且必须提供启用、验证、停用、卸载和回滚路径。

Compiler backends, kernel libraries, model conversion, benchmark suites, deployment managers,
profilers, and offline analysis tools are adjacent assets unless they register a vLLM runtime entry
point. Being used by a plugin does not make an adjacent asset a plugin.

编译器后端、算子库、模型转换工具、基准套件、部署管理器、性能分析器和离线分析工具属于相邻资产，
除非它们注册了 vLLM 运行时入口。被插件使用并不会使相邻资产自动成为插件。

\begin{normative}{Minimum boundary / 最小边界}
  A plugin \must be independently installable, selectable at process startup, identifiable in
  runtime evidence, and removable without editing the inference-engine source tree.

  插件必须能够独立安装、在进程启动时显式选择、在运行证据中可识别，并可在不修改推理引擎源码树
  的情况下卸载。
\end{normative}

\section{Supported entry-point groups / 支持的入口点组}

Choose the narrowest supported group that matches the extension's responsibility.

\begin{longtable}{>{\small\ttfamily\raggedright\arraybackslash}p{0.28\textwidth}>{\raggedright\arraybackslash}p{0.31\textwidth}>{\raggedright\arraybackslash}p{0.31\textwidth}}
\toprule
\normalfont\bfseries Group & \bfseries Intended use / 用途 & \bfseries Process scope / 进程范围 \\
\midrule
\endhead
vllm.general\_plugins & Runtime registration and general extensions / 通用运行时注册与扩展 & API, engine-core, worker \\
vllm.platform\_plugins & Out-of-tree hardware platform / 树外硬件平台 & Relevant initialization processes / 相关初始化进程 \\
vllm.io\_processor\_plugins & Input/output processors / 输入输出处理器 & API process \\
vllm.stat\_logger\_plugins & Statistics loggers / 统计日志器 & Asynchronous API process / 异步 API 进程 \\
\bottomrule
\end{longtable}

\texttt{vllm.platform\_plugins} is reserved for hardware-platform implementations. A compiler,
kernel library, model-preparation utility, or benchmark is not a platform plugin merely because a
runtime uses it.

\section{Package and identity / 包结构与身份}

\begin{lstlisting}
<plugin-repository>/
|-- pyproject.toml
|-- README.md
|-- CHANGELOG.md
|-- src/<package>/
|   |-- __init__.py
|   `-- plugin.py
`-- tests/
    |-- test_entry_point.py
    |-- test_registration.py
    `-- test_lifecycle.py
\end{lstlisting}

The package \must declare one stable plugin identity. Distribution name, import package, and plugin
ID may differ, but all three \must remain stable within a release line and be documented together.

\begin{lstlisting}
[project]
name = "<distribution-name>"
version = "<version>"

[project.entry-points."vllm.general_plugins"]
<plugin-id> = "<package>.plugin:register"
\end{lstlisting}

Every release \must record the distribution version, import package, plugin ID, entry-point group,
entry callable, supported vLLM range, and tested vLLM-HUST commit.

\section{Registration contract / 注册合同}

Registration \must be re-entrant because vLLM may load plugins in multiple processes and may invoke
the same registration path more than once.

\begin{lstlisting}[language=Python]
_REGISTERED = False

def register() -> None:
    global _REGISTERED
    if _REGISTERED:
        return
    # Install one narrow runtime hook.
    _REGISTERED = True
\end{lstlisting}

A conforming registration function:

\begin{itemize}
  \item \must install the narrowest practical hook and produce the same result on repeated calls;
  \item \must not duplicate patches, threads, sockets, files, or global state;
  \item \must validate required configuration before serving requests;
  \item \must emit one activation record containing plugin ID and version;
  \item \must fail visibly when a required hook or compatibility condition is unavailable;
  \item \must not embed machine paths, ports, device IDs, credentials, or model locations; and
  \item \should expose plugin-owned health, counters, or receipts when operationally relevant.
\end{itemize}

\section{Compatibility, configuration, and safety / 兼容、配置与安全}

Each release \must document:

\begin{itemize}
  \item supported platform, Python, vLLM, and vLLM-HUST versions;
  \item incompatible plugins and mutually exclusive runtime modes;
  \item every plugin-owned configuration variable, type, default, and validation rule;
  \item required external services or storage and their failure behavior; and
  \item safe fallback and rollback behavior.
\end{itemize}

Configuration variables \must use a plugin-specific prefix. Secrets \must not appear in source,
logs, manifests, command examples, or generated artifacts. A plugin \must fail closed when invalid
configuration could compromise correctness, isolation, or data integrity.

\section{Build, install, and discovery / 构建、安装与发现}

Build a wheel and install it into the exact Python environment used by every serving process.
Deployment-specific values remain external inputs.

\begin{lstlisting}[language=bash]
export BUILD_PYTHON=/path/to/build/python
export RUNTIME_PYTHON=/path/to/serving/python
export PLUGIN_WHEEL=/path/to/plugin.whl

"${BUILD_PYTHON}" -m build
"${RUNTIME_PYTHON}" -m pip install "${PLUGIN_WHEEL}"
\end{lstlisting}

Before startup, verify discovery using the same runtime interpreter:

\begin{lstlisting}[language=bash]
"${RUNTIME_PYTHON}" -c \
  'from importlib.metadata import entry_points; \
print([(e.name, e.value) for e in \
entry_points(group="vllm.general_plugins")])'
\end{lstlisting}

Installation alone does not prove activation. Multi-node deployments \must install and verify the
same plugin release in every applicable environment.

\section{Enable and start / 启用与启动}

Production services \must use an explicit allowlist.

\begin{lstlisting}[language=bash]
export PLUGIN_ID=your-plugin-id
export MODEL=your-model
export PORT=your-port

VLLM_PLUGINS="${PLUGIN_ID}" \
  vllm serve "${MODEL}" --port "${PORT}"
\end{lstlisting}

Multiple plugin IDs are comma-separated. If \texttt{VLLM\_PLUGINS} is unset, vLLM loads every
discovered plugin. If set to an empty string, vLLM loads none. Service-manager, container, model,
port, device, storage, and endpoint values \must remain deployment inputs.

\section{Verification and evidence / 验证与证据}

Activation requires all of the following:

\begin{enumerate}
  \item the intended entry point is discoverable in every serving environment;
  \item registration logs identify plugin ID and version without exposing secrets;
  \item the server reaches its readiness endpoint;
  \item at least one functional request completes correctly; and
  \item plugin-owned health, counters, or receipts are available when declared.
\end{enumerate}

\begin{lstlisting}[language=bash]
curl --fail --silent \
  "http://127.0.0.1:${PORT}/v1/models"
\end{lstlisting}

\begin{normative}{Evidence boundary / 证据边界}
  A smoke test validates discovery and lifecycle only. A performance claim \must use a matched
  baseline and treatment with the same model, request sequence, runtime, topology, and hardware
  allocation. Evidence \must identify whether it is real-online, trace-driven, replay, fixture,
  simulation, or projected; these labels are not interchangeable.

  Smoke 仅验证发现与生命周期。性能声明必须采用模型、请求序列、运行时、拓扑和硬件分配一致的
  matched 对照，并准确标记证据性质。
\end{normative}

\section{Stop and disable / 停止与停用}

Plugins are process-scoped. Standard 1.0 does not define hot unload.

\begin{enumerate}
  \item Gracefully stop the owning service through its supervisor or normal termination signal.
  \item Wait for API, engine-core, and worker processes to exit.
  \item Confirm plugin-owned threads, sockets, files, shared memory, and device resources are released.
  \item Remove the plugin ID from the explicit allowlist, or set \texttt{VLLM\_PLUGINS=""}.
  \item Restart and verify the baseline path.
\end{enumerate}

Do not use broad process-name termination. For a directly owned process whose PID was recorded at
startup, graceful termination is:

\begin{lstlisting}[language=bash]
kill -TERM "${SERVER_PID}"
wait "${SERVER_PID}"
\end{lstlisting}

For supervised deployments, use the supervisor's scoped stop operation rather than signaling an
unowned child process.

\section{Remove and roll back / 卸载与回滚}

Only uninstall after every consuming process has stopped:

\begin{lstlisting}[language=bash]
export PLUGIN_DISTRIBUTION=your-distribution-name
"${RUNTIME_PYTHON}" -m pip uninstall \
  "${PLUGIN_DISTRIBUTION}"
\end{lstlisting}

Then verify that the entry point is absent, restart the baseline with an explicit allowlist, run the
readiness probe, and retain the previous wheel and configuration until rollback verification is
complete. Rollback \must restore both software and plugin-owned state to a documented compatible
version.

\section{Conformance requirements / 符合性要求}

A release may be presented as \textbf{vLLM-HUST Plugin Standard 1.0 conforming} only when automated
tests cover every mandatory item below.

\begin{longtable}{>{\raggedright\arraybackslash}p{0.07\textwidth}>{\raggedright\arraybackslash}p{0.42\textwidth}>{\raggedright\arraybackslash}p{0.42\textwidth}}
\toprule
\textbf{ID} & \textbf{Required evidence} & \textbf{验收证据} \\
\midrule
\endhead
C01 & Wheel build and clean-environment installation & wheel 构建与干净环境安装 \\
C02 & Entry-point discovery in every required process environment & 各必要进程环境中的入口点发现 \\
C03 & Registration invoked twice without duplicate state & 重复注册不产生重复状态 \\
C04 & Registration in declared process roles & 在声明的进程角色中完成注册 \\
C05 & Enabled startup, readiness, and functional request & 启用后的启动、就绪与功能请求 \\
C06 & Disabled startup, readiness, and baseline request & 停用后的启动、就绪与基线路径 \\
C07 & Invalid configuration and compatibility rejection & 非法配置与不兼容条件拒绝 \\
C08 & Graceful shutdown and complete resource cleanup & 优雅停止与完整资源清理 \\
C09 & Baseline restart after disable or uninstall & 停用或卸载后的基线重启 \\
C10 & Versioned evidence manifest with software and configuration identity & 含软件和配置身份的版本化证据 manifest \\
\bottomrule
\end{longtable}

The plugin README \must contain exact build, install, discover, enable, start, verify, stop, disable,
uninstall, and rollback instructions for the released distribution. A conformance statement
\should link to the test run and immutable release artifact.

\section{Release checklist / 发布检查清单}

\begin{tcolorbox}[colback=white,colframe=Teal,arc=0pt,title=\textbf{Operator handoff / 运维交付}]
\begin{itemize}[label=$\square$]
  \item Stable package identity and entry point / 稳定的包身份与入口点
  \item Compatibility matrix and configuration reference / 兼容矩阵与配置说明
  \item Signed or hashed release artifact / 签名或哈希固定的发布产物
  \item Explicit enable and disable commands / 明确的启用与停用命令
  \item Readiness, correctness, and cleanup evidence / 就绪、正确性与清理证据
  \item Rollback artifact and procedure / 回滚产物与操作流程
  \item No credentials or deployment-specific constants / 不含凭据或部署专用常量
\end{itemize}
\end{tcolorbox}

\section{References / 参考资料}

\begin{itemize}
  \item \href{https://docs.vllm.ai/en/latest/design/plugin_system/}{vLLM Plugin System}
  \item \href{https://packaging.python.org/en/latest/specifications/entry-points/}{Python Packaging Entry Points Specification}
  \item \href{https://vllm-hust.sage.org.ai/plugins.html}{vLLM-HUST Plugin Ecosystem}
\end{itemize}

\vfill
\begin{center}
\textcolor{Slate}{\small End of vLLM-HUST Plugin Standard 1.0 / 标准正文结束}
\end{center}

\end{document}
